\clearpage\thispagestyle{empty}
\pdfbookmark[0]{Quellen zu Gauß an von Lindenau, 30. November 1808}{lindenau1808nov-sources}
\begin{center}\textbf{Quellen zu Gauß an von Lindenau, 30. November 1808}\end{center}
\noindent Diese KI-Übertragung in heutiges Deutsch folgt Carl Friedrich Gauß’ deutschem Brief in den \textit{Werken}, Band~VI (Göttingen, 1874), S.~141--145, einschließlich der handschriftlichen Bemerkung. Der \href{https://archive.org/details/bub_gb_caMAAAAAMAAJ/page/n91/mode/1up}{Erstdruck}, Beitrag V im Januarheft 1809 der von Franz Xaver von Zach herausgegebenen \textit{Monatlichen Correspondenz zur Beförderung der Erd- und Himmels-Kunde}, Band~19, S.~85--92, enthält auf S.~86 eine mit v.~L. unterzeichnete Anmerkung und auf S.~90 Gauß’ Liste der Druckfehler in seiner Abhandlung vom Oktober. \href{https://gdz.sub.uni-goettingen.de/id/PPN236007467}{Historischer deutscher Text in den Werken}.
\par\smallskip\noindent Joseph-Jérôme de Lalande beschreibt Cagnolis Verfahren und das Rechenbeispiel in seiner \href{https://doi.org/10.3931/e-rara-9241}{\textit{Astronomie}}, 3.~Aufl., Band~III (Paris, 1792), Art.~3259--3261. Die von Gauß erwähnten Logarithmentafeln sind Michael Taylors \href{https://doi.org/10.3931/e-rara-9132}{\textit{Tables of Logarithms}} (London, 1792), mit Vorwort und Anleitungen von Nevil Maskelyne. Franz Brünnow behandelt Gauß’ Verfahren mit einem eigenen Beispiel im \href{https://archive.org/details/bub_gb_sAA5AAAAMAAJ/page/n306/mode/1up}{\textit{Lehrbuch der sphärischen Astronomie}} (Berlin, 1851), §20, und in seiner englischen Übersetzung der zweiten deutschen Auflage, \href{https://archive.org/details/sphericalastro00brhurich/page/n328/mode/1up}{\textit{Spherical Astronomy}} (London, 1865), §20.
\par\smallskip\noindent Übertragung und Satz wurden mit KI erstellt.
\clearpage
